\documentclass[12pt]{article}
%\usepackage[sort&compress,comma,numbers,round]{natbib}
\usepackage{graphicx}
\usepackage{bm}
\usepackage[hidelinks=true]{hyperref} 
\usepackage{xcolor}
\usepackage{amsmath}
\usepackage{amsfonts}
\usepackage{float, bigints}
\usepackage{amssymb,graphicx,fullpage}
\usepackage{mathrsfs}
\usepackage{bm}
\usepackage[hidelinks=true]{hyperref} 
\usepackage{mathrsfs, amsmath}
\usepackage{amsbsy,setspace}
\usepackage{authblk}

\usepackage[most]{tcolorbox}
\usepackage{comment}
\tcbuselibrary{theorems}

\newtcbtheorem[]{mytheo}{}%
{colback=gray!5,colframe=gray!35!black,fonttitle=\bfseries}{th}

\renewcommand\Affilfont{\itshape\small}
%\definecolor{greenish}{rgb}{0.13,0.58,0.16}
\definecolor{greenish}{RGB}{141,198,63}
\definecolor{reddish}{RGB}{239,65,54}
\definecolor{blueish}{RGB}{28,117,188}

\hypersetup{
colorlinks   = true, %Colours links instead of ugly boxes
urlcolor     = blueish, %Colour for external hyperlinks
linkcolor    = magenta, %Colour of internal links
citecolor   =  greenish%Colour of citations
}

\def\d{\mathrm{d}}
\def\e{\mathrm{e}}
\def\hp{\hat{p}}
\def\E{\mathbb{E}}
\def\eh{\hat{\mathbf{e}}}
\def\et{\tilde{\mathbf{e}}}
\def\a{\mathbf{a}}
\def\b{\mathbf{b}}
\def\c{\mathbf{c}}
\def\v{\mathbf{v}}
\def\p{\mathbf{p}}
\def\q{\mathbf{q}}
\def\p{\mathbf{a}}
\def\F{\mathbb{F}}
\def\B{\mathbb{B}}
\def\A{\mathbf{A}}
\def\C{\mathbb{C}}
\def\R{\mathbb{R}}
\def\J{\mathcal{J}}
\def\T{\dagger}
\def\fh{\hat{f}}
\def\N{\mathcal{N}}
\def\var{\mathrm{Var}}
\def\Z{\mathbb{Z}}
\def\grad{\nabla}
\def\P{\mathbb{P}}
\def\D{\mathbb{D}}
\def\phih{\hat{\phi}}
\def\alpha{R}
\def\q{Q}
\def\ss{\mathrm{ss}}
\def\sgp#1{\textcolor{red}{#1}}

\title{Probability, Stochastic Processes and Path Integrals}
\author{S Ganga Prasath\thanks{\href{mailto:sgangaprasath@smail.iitm.ac.in}{sgangaprasath@smail.iitm.ac.in}}}
\date{}

\begin{document}

\maketitle
\begin{singlespace}

\section{Classic distributions}
\begin{itemize}
    \item \textbf{Gaussian distribution}: If $X \sim \mathcal{N}(\mu, \sigma^2)$, then the probability density function (pdf) of $X$ is given by:
    \[
    f_X(x) = \frac{1}{\sqrt{2\pi\sigma^2}} \exp\left(-\frac{(x - \mu)^2}{2\sigma^2}\right)
    \]
    The Gaussian distribution is symmetric around its mean $\mu$ and has a bell-shaped curve. It is widely used in statistics and machine learning due to the Central Limit Theorem.

    \item \textbf{Bernoulli distribution}: If $X \sim \text{Bernoulli}(p)$, then the probability mass function (pmf) of $X$ is given by:
    \[
    P(X = x) = p^x (1-p)^{1-x} \quad \text{for } x = 0, 1
    \]
    The Bernoulli distribution models a binary outcome, where $p$ is the probability of success (i.e., $X=1$).

    \item \textbf{Binomial distribution}: If $X \sim \text{Binomial}(n, p)$, then the pmf of $X$ is given by:
    \[
    P(X = k) = \binom{n}{k} p^k (1-p)^{n-k} \quad \text{for } k = 0, 1, ..., n
    \]
    The Binomial distribution models the number of successes in $n$ independent Bernoulli trials with success probability $p$.

    \item \textbf{Poisson distribution}: If $X \sim \text{Poisson}(\lambda)$, then the pmf of $X$ is given by:
    \[
    P(X = k) = e^{-\lambda} \frac{\lambda^k}{k!} \quad \text{for } k = 0, 1, 2, ...
    \]
    The Poisson distribution models the number of events occurring in a fixed interval of time or space when events occur with a known constant mean rate $\lambda$ and independently of the time since the last event.
\end{itemize}
\section{Useful properties of distributions and random variables}
\begin{itemize}
    \item \textbf{Moment generating function:} The moment generating function (MGF) of a random variable $X$ is defined as $M_X(t) = \mathbb{E}[e^{tX}]$. It can be used to derive moments and other properties of the distribution. Here $t$ is a real number, and the MGF exists for values of $t$ in an interval around 0 where the expectation is finite. The n-th moment of $X$ can be obtained by taking the n-th derivative of the MGF at $t=0$:
    \[
    \mathbb{E}[X^n] = M_X^{(n)}(0).
    \]
    \item Example: For a Gaussian random variable $X \sim \mathcal{N}(\mu, \sigma^2)$, the MGF is given by:
    \[
    M_X(t) = \exp\left(\mu t + \frac{1}{2} \sigma^2 t^2\right).
    \]
    We can derive this expression by evaluating the expectation:
    \[
    M_X(t) = \mathbb{E}[e^{tX}] = \int_{-\infty}^{\infty} e^{tx} f_X(x) \,dx = \int_{-\infty}^{\infty} e^{tx} \frac{1}{\sqrt{2\pi\sigma^2}} \exp\left(-\frac{(x - \mu)^2}{2\sigma^2}\right) dx.
    \]
    The MGF can be used to derive the mean and variance of $X$ by taking derivatives:
    \[
    \mathbb{E}[X] = M_X'(0) = \mu, \quad \text{Var}(X) = M_X''(0) - (M_X'(0))^2 = \sigma^2.
    \]
    For a Bernoulli random variable $X \sim \text{Bernoulli}(p)$, the MGF is given by:
    \[  M_X(t) = \mathbb{E}[e^{tX}] = p e^t + (1-p) e^0 = p e^t + (1-p). \]
    We can derive the mean and variance of $X$ using the MGF:
    \[
    \mathbb{E}[X] = M_X'(0) = p, \quad \text{Var}(X) = M_X''(0) - (M_X'(0))^2 = p(1-p).
    \]
    \item \textbf{Cumulant generating function:} The cumulant generating function (CGF) of a random variable $X$ is defined as $K_X(t) = \log M_X(t)$. The CGF can be used to derive cumulants, which are related to moments but have properties that make them useful for certain applications, such as in the study of distributions and their convergence. The cumulant is related to the moments of a random variable, and the $n$-th cumulant can be obtained by taking the $n$-th derivative of the CGF at $t=0$:
    \[
    \kappa_n = K_X^{(n)}(0).
    \] The first cumulant is the mean, the second cumulant is the variance, and higher-order cumulants provide information about the shape of the distribution (e.g., skewness and kurtosis).\\
    The reason for using cumulants instead of moments is that cumulants have properties that make them more suitable for certain applications, such as in the study of distributions and their convergence. Further, the reason $K_X(t) = \log M_X(t)$ is that the logarithm of the MGF transforms the product of MGFs (which corresponds to the sum of independent random variables) into a sum of CGFs, making it easier to analyze the behavior of sums of random variables. Additionally, the CGF can provide insights into the tail behavior and other properties of the distribution that may not be as easily accessible through moments alone.
    \item  Example: For a Gaussian random variable $X \sim \mathcal{N}(\mu, \sigma^2)$, the CGF is given by:
    \[
    K_X(t) = \log M_X(t) = \mu t + \frac{1}{2} \sigma^2 t^2.
    \] The first cumulant is $\kappa_1 = K_X'(0) = \mu$, and the second cumulant is $\kappa_2 = K_X''(0) = \sigma^2$. For a Bernoulli random variable $X \sim \text{Bernoulli}(p)$, the CGF is given by:
    \[
    K_X(t) = \log M_X(t) = \log(p e^t + (1-p)).
    \] The first cumulant is $\kappa_1 = K_X'(0) = p$, and the second cumulant is $\kappa_2 = K_X''(0) = p(1-p)$.
    \item \textbf{Characteristic function:} The characteristic function of a random variable $X$ is defined as $\phi_X(t) = \mathbb{E}[e^{itX}]$. It can be used to derive moments and other properties of the distribution, and it is particularly useful for studying the distribution of sums of independent random variables.
    \item The characteristic function and the moment generating function are related by $M_X(t) = \phi_X(-it)$.
\end{itemize}

\section{Stable distributions}
A stable distribution is a type of probability distribution that has the property of stability under convolution. This means that if you take two independent random variables that follow a stable distribution and add them together, the resulting random variable will also follow a stable distribution (up to location and scale parameters). Stable distributions are characterized by four parameters: the stability parameter $\alpha$, the skewness parameter $\beta$, the scale parameter $\gamma$, and the location parameter $\delta$. The stability parameter $\alpha$ determines the tail behavior of the distribution, with smaller values indicating heavier tails. The skewness parameter $\beta$ controls the asymmetry of the distribution, while the scale parameter $\gamma$ affects the spread of the distribution, and the location parameter $\delta$ shifts the distribution along the real line. Stable distributions include well-known distributions such as the Gaussian distribution (when $\alpha=2$ and $\beta=0$), Levy distribution (when $\alpha=1.5$ and $\beta=0$) and the Cauchy distribution (when $\alpha=1$ and $\beta=0$). They are used in various fields, including finance, physics, and signal processing, to model phenomena with heavy-tailed behavior or skewed distributions.

\subsection{Gaussian distribution}
Let us consider a random variable $X$ that follows Gaussian distribution with mean $\mu$ and variance $\sigma^2$. The probability density function (pdf) of $X$ is given by:
\[
f_X(x) = \frac{1}{\sqrt{2\pi\sigma^2}} \exp\left(-\frac{(x - \mu)^2}{2\sigma^2}\right).
\]
Now, consider sum of $n$ independent and identically distributed (i.i.d.) random variables $X_1, X_2, ..., X_n$, each following the same Gaussian distribution as $X$. The sum of these random variables is given by:
\[
S_n = X_1 + X_2 + ... + X_n.
\]
The sum of $n$ independent Gaussian random variables also follows a Gaussian distribution, with mean $n\mu$ and variance $n\sigma^2$. We can derive this result using the properties of the moment generating function (MGF). The MGF of a Gaussian random variable $X$ is given by:
\[
M_X(t) = \exp\left(\mu t + \frac{1}{2} \sigma^2 t^2\right).
\]
Since $X_1, X_2, ..., X_n$ are independent and identically distributed, the MGF of their sum $S_n$ is the product of their individual MGFs:
\[
M_{S_n}(t) = [M_X(t)]^n = \exp\left(n\mu t + \frac{1}{2} n\sigma^2 t^2\right).
\]
This is the MGF of a Gaussian distribution with mean $n\mu$ and variance $n\sigma^2$. Therefore, we can conclude that the sum of $n$ independent Gaussian random variables follows a Gaussian distribution with mean $n\mu$ and variance $n\sigma^2$.

\subsection{Levy distribution}
The Levy distribution, denoted as $L(\delta, \gamma)$, is a continuous probability distribution with several interesting properties. For a random variable $Y \sim \text{Levy}(\delta, \gamma)$, the PDF is:
\begin{equation}
f(y; \delta, \gamma) = \sqrt{\frac{\gamma}{2\pi}} \frac{\exp\left( -\frac{\gamma}{2(y-\delta)} \right)}{(y-\delta)^{3/2}}, \quad y > \delta
\end{equation}
where $\delta$ is the location parameter and $\gamma$ is the scale parameter.
The CDF is expressed in terms of the complementary error function ($\text{erfc}$):
\begin{equation}
F(y; \delta, \gamma) = \text{erfc}\left( \sqrt{\frac{\gamma}{2(y-\delta)}} \right)
\end{equation}
The Levy distribution is characterized by ``heavy tails,'' specifically:
\begin{itemize}
    \item \textbf{Mean:} $\mathbb{E}[Y] = \infty$. The integral for the first moment diverges.
    \item \textbf{Variance:} $\text{Var}(Y) = \infty$.
    \item \textbf{Tail Behavior:} As $y \to \infty$, the PDF decays as $f(y) \sim y^{-3/2}$. This power-law decay is much slower than the exponential decay of a Normal distribution.
\end{itemize}
The Levy distribution is strictly stable. If $Y_1, Y_2, \dots, Y_n$ are i.i.d. Levy random variables with parameters $(\delta, \gamma)$, then their sum $S_n$ follows:
\begin{equation}
S_n = \sum_{j=1}^n Y_j \sim \text{Levy}(n\delta, n\gamma)
\end{equation}
This implies that the shape of the distribution is preserved under convolution, provided the location and scale parameters are adjusted.

Let us now consider a random variable $Y$ that follows a Levy distribution with location parameter $\delta$ and scale parameter $\gamma$. The probability density function (pdf) of $Y$ is given by:
\[
f_Y(y) = \sqrt{\frac{\gamma}{2\pi}} \frac{\exp\left(-\frac{\gamma}{2(y-\delta)}\right)}{(y-\delta)^{3/2}} \quad \text{for } y > \delta.
\]
Consider the sum of $n$ independent and identically distributed (i.i.d.) random variables $Y_1, Y_2, ..., Y_n$, each following the same Levy distribution as $Y$. The sum of these random variables is given by:
\[
S_n = Y_1 + Y_2 + ... + Y_n.
\]
The sum of $n$ independent Levy random variables also follows a Levy distribution, with location parameter $n\delta$ and scale parameter $n\gamma$. We can derive this result using the properties of the characteristic function. The characteristic function of a Levy random variable $Y$ can be calculated as follows:
\[
\phi_Y(t) = \mathbb{E}[e^{itY}] = \int_{\delta}^{\infty} e^{ity} f_Y(y) \,dy = \int_{\delta}^{\infty} e^{ity} \sqrt{\frac{\gamma}{2\pi}} \frac{\exp\left(-\frac{\gamma}{2(y-\delta)}\right)}{(y-\delta)^{3/2}} \,dy.
\]
We can simplify this expression by making the substitution $u = y - \delta$, which gives us:
\[
\phi_Y(t) = e^{i\delta t} \int_{0}^{\infty} e^{itu} \sqrt{\frac{\gamma}{2\pi}} \frac{\exp\left(-\frac{\gamma}{2u}\right)}{u^{3/2}} \,du.
\]
 We can simplify this expression further to get:
\[
\int_{0}^{\infty} e^{itu} \sqrt{\frac{\gamma}{2\pi}} \frac{\exp\left(-\frac{\gamma}{2u}\right)}{u^{3/2}} \,du = \sqrt{\frac{\gamma}{2\pi}} \int_{0}^{\infty} u^{-3/2} \exp\left(-\frac{\gamma}{2u} + itu\right) \,du.
\]
To evaluate the integral:
\begin{equation}
I = \int_{0}^{\infty} u^{-3/2} \exp\left(-\frac{a}{u} - bu\right) du, \quad a, b > 0,
\end{equation}
let us substitute $u = x^2$, which implies $du = 2x \, dx$. Substituting these into the integral:
\begin{align}
I &= \int_{0}^{\infty} (x^2)^{-3/2} \exp\left(-\frac{a}{x^2} - bx^2\right) (2x) \, dx = 2 \int_{0}^{\infty} \frac{1}{x^2} \exp\left(-\frac{a}{x^2} - bx^2\right) dx.
\end{align}
We rewrite the argument of the exponential by completing the square:
\begin{equation}
-\left(\frac{a}{x^2} + bx^2\right) = -\left(\frac{\sqrt{a}}{x} - x\sqrt{b}\right)^2 - 2\sqrt{ab}.
\end{equation}
Substituting this back into the integral:
\begin{equation}
I = 2e^{-2\sqrt{ab}} \int_{0}^{\infty} \frac{1}{x^2} \exp\left[ -\left(\frac{\sqrt{a}}{x} - x\sqrt{b}\right)^2 \right] dx.
\end{equation}
We wish to evaluate the integral:
\begin{equation}
J = \int_{0}^{\infty} \frac{1}{x^2} \exp\left[ -\left(x\sqrt{b} - \frac{\sqrt{a}}{x}\right)^2 \right] dx
\end{equation}
Let $z = \frac{\sqrt{a}}{x}$, which implies $x = \frac{\sqrt{a}}{z}$ and $dx = -\frac{\sqrt{a}}{z^2} dz$.
As $x \to 0$, $z \to \infty$; as $x \to \infty$, $z \to 0$. Substituting these into $J$:
\begin{equation}
J = \frac{1}{\sqrt{a}} \int_{0}^{\infty} \exp\left[ -\left(z - \frac{\sqrt{ab}}{z}\right)^2 \right] dz
\end{equation}
Note that the exponent is symmetric. We use the identity for any integrable function $f$:
\begin{equation}
\int_{0}^{\infty} f\left( z - \frac{k}{z} \right) dz = \int_{0}^{\infty} f(v) dv \quad \text{where } k = \sqrt{ab}
\end{equation}
Applying this to our integral where $f(v) = e^{-v^2}$:
\begin{equation}
\int_{0}^{\infty} \exp\left[ -\left(z - \frac{\sqrt{ab}}{z}\right)^2 \right] dz = \int_{0}^{\infty} e^{-v^2} dv
\end{equation}
We know the standard Gaussian integral $\int_{0}^{\infty} e^{-v^2} dv = \frac{\sqrt{\pi}}{2}$.

Substituting the Gaussian result back into our expression for $J$:
\begin{equation}
J = \frac{1}{\sqrt{a}} \left( \frac{\sqrt{\pi}}{2} \right) = \sqrt{\frac{\pi}{4a}}
\end{equation}

Now substitute $J$ back into the full expression for $I$ from Step 2:
\begin{align}
I &= 2e^{-2\sqrt{ab}} J \\
I &= 2e^{-2\sqrt{ab}} \sqrt{\frac{\pi}{4a}} \\
I &= \sqrt{\frac{\pi}{a}} e^{-2\sqrt{ab}}
\end{align}

Let $z = x\sqrt{b} - \frac{\sqrt{a}}{x}$. To handle the differential, we use the substitution $y = \frac{\sqrt{a}}{x\sqrt{b}}$, which shows that:
\begin{equation}
\int_{0}^{\infty} \frac{1}{x^2} \exp\left[ -\left(\frac{\sqrt{a}}{x} - x\sqrt{b}\right)^2 \right] dx = \sqrt{\frac{\pi}{4a}}
\end{equation}
More formally, using the identity for integrals of the form $f(x\sqrt{b} - \sqrt{a}/x)$, the integral evaluates as:
\begin{equation}
I = 2e^{-2\sqrt{ab}} \left( \frac{1}{2} \sqrt{\frac{\pi}{a}} \right)
\end{equation}
The closed-form solution is:
\begin{equation}
\int_{0}^{\infty} u^{-3/2} \exp\left(-\frac{a}{u} - bu\right) du = \sqrt{\frac{\pi}{a}} e^{-2\sqrt{ab}}
\end{equation}
Given the constant $\sqrt{\frac{\gamma}{2\pi}}$ and parameters $a = \frac{\gamma}{2}, b = -it$:
\begin{align}
\phi_Y(t) &= e^{i\delta t} \sqrt{\frac{\gamma}{2\pi}} \left( \sqrt{\frac{\pi}{\gamma/2}} \exp\left(-2\sqrt{\frac{\gamma}{2}(-it)}\right) \right) \\
\phi_Y(t) &= e^{i\delta t} \sqrt{\frac{\gamma}{2\pi}} \sqrt{\frac{2\pi}{\gamma}} \exp\left(-\sqrt{-2i\gamma t}\right) \\
\phi_Y(t) &= \exp\left(i\delta t - \sqrt{-2i\gamma t}\right)
\end{align}
Ultimately we get the characteristic function of a Levy random variable as:
\[
\phi_Y(t) = \exp\left(i\delta t - \gamma |t|^{1/2}\right).
\]
Since $Y_1, Y_2, ..., Y_n$ are independent and identically distributed, the characteristic function of their sum $S_n$ is the product of their individual characteristic functions:
\[
\phi_{S_n}(t) = [\phi_Y(t)]^n = \exp\left(i n\delta t - n\gamma |t|^{1/2}\right).
\]
This is the characteristic function of a Levy distribution with location parameter $n\delta$ and scale parameter $n\gamma$. Therefore, we can conclude that the sum of $n$ independent Levy random variables follows a Levy distribution with location parameter $n\delta$ and scale parameter $n\gamma$.

\section{Central Limit Theorem}
The Central Limit Theorem (CLT) is a fundamental result in probability theory that states that the sum of a large number of independent and identically distributed (i.i.d.) random variables, each with finite mean and variance, will approximately follow a normal distribution, regardless of the original distribution of the variables. 

Consider i.i.d. Bernoulli variables $X_i$ with mean $\mu = p$ and variance $\sigma^2 = p(1-p)$. To reach the standard normal distribution, we must use the standardized variable:
\begin{equation}
Y_i = \frac{X_i - p}{\sigma} = \frac{X_i - p}{\sqrt{p(1-p)}}
\end{equation}
The charasteristic function of $X_i$ is given by:
\[
\phi_X(t) = \mathbb{E}[e^{itX_i}] = p e^{it} + (1-p) e^{i0} = p e^{it} + (1-p).
\]
The characteristic function of $Y_i$, denoted by $\phi_{Y}(t)$, is related to the characteristic function of $X_i$ by the properties of linear transformations $\phi_{aX+b}(t) = e^{itb}\phi_X(at)$:
\begin{equation}
\phi_{Y}(t) = e^{-i\frac{p}{\sigma}t} \phi_{X}\left(\frac{t}{\sigma}\right) = e^{-i\frac{p}{\sigma}t} \left( p e^{i\frac{t}{\sigma}} + (1-p) \right)
\end{equation}

We expand the terms inside the parentheses and the exponential for small $t$:
\begin{align}
\phi_{Y}(t) &= \left( 1 - \frac{ipt}{\sigma} - \frac{p^2 t^2}{2\sigma^2} + \dots \right) \left( p \left[ 1 + \frac{it}{\sigma} - \frac{t^2}{2\sigma^2} + \dots \right] + 1 - p \right) \\
&= \left( 1 - \frac{ipt}{\sigma} - \frac{p^2 t^2}{2\sigma^2} \right) \left( 1 + \frac{ipt}{\sigma} - \frac{pt^2}{2\sigma^2} \right) + O(t^3)
\end{align}
Multiplying these terms and keeping only up to $t^2$:
\begin{equation}
\phi_{Y}(t) = 1 + \left( \frac{ipt}{\sigma} - \frac{ipt}{\sigma} \right) + \left( \frac{p^2 t^2}{\sigma^2} - \frac{pt^2}{2\sigma^2} - \frac{p^2 t^2}{2\sigma^2} \right) = 1 - \frac{p(1-p)t^2}{2\sigma^2}
\end{equation}
Since $\sigma^2 = p(1-p)$, the $p$ terms cancel perfectly:
\begin{equation}
\phi_{Y}(t) = 1 - \frac{t^2}{2} + O(t^3)
\end{equation}

Let $Y_1, Y_2, \dots, Y_n$ be i.i.d. random variables and we are interested in the characteristic function of $Z_n = \frac{1}{\sqrt{n}} \sum Y_i$. This is given by:
\begin{equation}
\phi_{Z_n}(t) = \left[ \phi_Y \left( \frac{t}{\sqrt{n}} \right) \right]^n = \left[ 1 - \frac{(t/\sqrt{n})^2}{2} + O(n^{-3/2}) \right]^n
\end{equation}
\begin{equation}
\phi_{Z_n}(t) = \left[ 1 - \frac{t^2}{2n} + O(n^{-3/2}) \right]^n
\end{equation}
Taking $n \to \infty$ and using the identity $\lim_{n \to \infty} (1 + x/n)^n = e^x$, we get:
\begin{equation}
\lim_{n \to \infty} \phi_{Z_n}(t) = \exp\left( -\frac{t^2}{2} \right)
\end{equation}
This is the characteristic function of the standard normal distribution $\mathcal{N}(0, 1)$. Given that the characteristic function uniquely determines the distribution, we conclude that $Z_n$ converges in distribution to $\mathcal{N}(0, 1)$ as $n \to \infty$. This is the essence of the Central Limit Theorem.

\section{Generalized Central Limit Theorem}
The Generalized Central Limit Theorem (GCLT) extends the classic Central Limit Theorem to cases where the underlying distribution of the random variables has infinite variance. Specifically, it states that if we have a sequence of independent and identically distributed (i.i.d.) random variables with heavy-tailed distributions (i.e., distributions that decay slower than a Gaussian), then the properly normalized sum of these random variables converges in distribution to a stable law, which may not be Gaussian.

Let $Y_1, Y_2, \dots, Y_n$ be i.i.d. Levy random variables with location $\delta = 0$ and scale $\gamma$. We have already derived the characteristic function:
\begin{equation}
\phi_Y(t) = \exp\left( -\sqrt{-2i\gamma t} \right)
\end{equation}
We define the normalized sum $Z_n = \frac{1}{n^2} \sum_{j=1}^n Y_j$. The characteristic function of $Z_n$ is:
\begin{align}
\phi_{Z_n}(t) &= \left[ \phi_Y\left( \frac{t}{n^2} \right) \right]^n \\
&= \left[ \exp\left( -\sqrt{-2i\gamma \frac{t}{n^2}} \right) \right]^n \\
&= \exp\left( -n \frac{\sqrt{-2i\gamma t}}{n} \right) \\
&= \exp\left( -\sqrt{-2i\gamma t} \right)
\end{align}
Thus, for any $n$, $\phi_{Z_n}(t) = \phi_{Z_\infty}(t) = \exp\left( -\sqrt{-2i\gamma t} \right)$.

To find the PDF $f_{Z_\infty}(z)$, we use the inversion formula:
\begin{equation}
f_{Z_\infty}(z) = \frac{1}{2\pi} \int_{-\infty}^{\infty} e^{-itz} \phi_{Z_\infty}(t) dt = \frac{1}{2\pi} \int_{-\infty}^{\infty} e^{-itz} e^{-\sqrt{-2i\gamma t}} dt
\end{equation}

This integral can be solved using a contour in the complex plane (Bromwich integral) or by recognizing it as a known Laplace transform pair. Let $s = -it$, then the integral takes the form of the inverse Laplace transform:
\begin{equation}
f_{Z_\infty}(z) = \mathcal{L}^{-1} \left\{ e^{-\sqrt{2\gamma s}} \right\}
\end{equation}
Using the standard transform pair $\mathcal{L}^{-1} \{ e^{-k\sqrt{s}} \} = \frac{k}{2\sqrt{\pi z^3}} \exp\left(-\frac{k^2}{4z}\right)$ with $k = \sqrt{2\gamma}$:
\begin{align}
f_{Z_\infty}(z) &= \frac{\sqrt{2\gamma}}{2\sqrt{\pi z^3}} \exp\left( -\frac{(\sqrt{2\gamma})^2}{4z} \right) \\
f_{Z_\infty}(z) &= \frac{\sqrt{2\gamma}}{2\sqrt{\pi} z^{3/2}} \exp\left( -\frac{2\gamma}{4z} \right)
\end{align}
Simplifying the constants:
\begin{equation}
f_{Z_\infty}(z) = \sqrt{\frac{\gamma}{2\pi}} \frac{1}{z^{3/2}} \exp\left( -\frac{\gamma}{2z} \right), \quad z > 0
\end{equation}
This confirms that the limit distribution $Z_\infty$ is identical in form to the original Levy distribution, proving that the Levy distribution is the ``attractor'' for sums of variables with $\alpha=1/2$.

\section{Path integral}
In this section we introduce the concept of path integrals, which are a powerful tool in both physics and probability theory for analyzing systems by integrating over all possible paths or trajectories that the system can take. The path integral formulation allows us to compute probabilities and expectations by summing over all possible configurations of a system, weighted by an appropriate action or energy functional.

Let us consider a stochastic process $X(t)$ defined on the interval $[0, T]$. The path integral formulation allows us to express the probability of a particular trajectory or path of the process as an integral over all possible paths. For example, if we want to compute the probability that the process follows a specific path $\gamma(t)$, we can write:
\begin{equation}
P(\gamma) = \int \mathcal{D}[X] \, e^{-S[X]} \delta(X(t) - \gamma(t))
\end{equation}
where $\mathcal{D}[X]$ denotes the integration over all paths $X(t)$, $S[X]$ is the action functional that assigns a weight to each path based on its properties, and $\delta$ is the Dirac delta function that ensures we are only considering paths that match the specified trajectory $\gamma(t)$. The action functional $S[X]$ typically encodes the dynamics of the system and can be derived from the underlying stochastic differential equations governing the process. The path integral formulation is particularly useful for analyzing systems with complex dynamics, such as those with non-linear interactions or those that exhibit phase transitions. It provides a framework for computing correlation functions, transition probabilities, and other quantities of interest by integrating over the space of all possible paths, allowing us to capture the full range of behaviors exhibited by the system.

\subsection{Gaussian process}
Before we delve into the path integral formulation for a general stochastic process, let us first consider the case of a Gaussian process. A Gaussian process is a collection of random variables, any finite number of which have a joint Gaussian distribution. It is fully characterized by its mean function $\mu(t)$ and covariance function $C(t, t')$. The probability density functional for a Gaussian process can be expressed as:
\begin{equation}P[X] \propto \exp\left( -\frac{1}{2} \int_0^T dt \int_0^T dt' (X(t) - \mu(t)) C^{-1}(t, t') (X(t') - \mu(t')) \right)\end{equation}
where $C^{-1}(t, t')$ is the inverse of the covariance function. This expression can be derived from the properties of Gaussian distributions and serves as a starting point for understanding how to construct the action functional for more general stochastic processes.

To derive the probability distribution for correlated variables from scratch, we start with the simplest case (independent variables) and use a \textbf{linear transformation} to ``mix'' them. This is the most intuitive way to see where the covariance matrix $C$ and its inverse $C^{-1}$ actually come from.
Suppose we have a set of $n$ independent Gaussian random variables $Z_1, Z_2, \dots, Z_n$, each with mean $0$ and variance $1$. Their joint probability is simply the product of $n$ standard normals:
\begin{equation}
P(\mathbf{z}) = \prod_{i=1}^n \frac{1}{\sqrt{2\pi}} \exp\left( -\frac{z_i^2}{2} \right) = \frac{1}{(2\pi)^{n/2}} \exp\left( -\frac{1}{2} \mathbf{z}^T \mathbf{z} \right)
\end{equation}
Here, the variables are uncorrelated. The ``covariance matrix'' is just the identity matrix $\mathbf{I}$. To create \textbf{correlated} variables $X$, we define a linear transformation:
\begin{equation}
\mathbf{x} = \mathbf{A}\mathbf{z} + \boldsymbol{\mu}
\end{equation}
where $\mathbf{A}$ is an $n \times n$ matrix that ``mixes'' the independent $z$'s together, and $\boldsymbol{\mu}$ is the mean vector.

To find the probability density $P(\mathbf{x})$, we use the multidimensional change of variables formula:
\begin{equation}
P(\mathbf{x}) = P(\mathbf{z}) \left| \det\left( \frac{d\mathbf{z}}{d\mathbf{x}} \right) \right|
\end{equation}
From our transformation, $\mathbf{z} = \mathbf{A}^{-1}(\mathbf{x} - \boldsymbol{\mu})$. The Jacobian is $\frac{d\mathbf{z}}{d\mathbf{x}} = \mathbf{A}^{-1}$. Thus:
\begin{equation}
P(\mathbf{x}) = \frac{1}{(2\pi)^{n/2} |\det(\mathbf{A})|} \exp\left( -\frac{1}{2} \left[ \mathbf{A}^{-1}(\mathbf{x} - \boldsymbol{\mu}) \right]^T \left[ \mathbf{A}^{-1}(\mathbf{x} - \boldsymbol{\mu}) \right] \right)
\end{equation}
The covariance of $\mathbf{x}$ can be written as:
\begin{equation}
\mathbf{C} = \mathbb{E}[(\mathbf{x}-\boldsymbol{\mu})(\mathbf{x}-\boldsymbol{\mu})^T] = \mathbb{E}[(\mathbf{A}\mathbf{z})(\mathbf{A}\mathbf{z})^T] = \mathbf{A} \mathbb{E}[\mathbf{z}\mathbf{z}^T] \mathbf{A}^T
\end{equation}
Since $Z$ are independent with variance $1$, $\mathbb{E}[\mathbf{z}\mathbf{z}^T] = \mathbf{I}$. Therefore:
\begin{equation}
\mathbf{C} = \mathbf{A}\mathbf{A}^T
\end{equation}
This confirms that the term in our exponent $(\mathbf{A}\mathbf{A}^T)^{-1}$ is exactly $\mathbf{C}^{-1}$. Also, note that $\det(\mathbf{C}) = \det(\mathbf{A}\mathbf{A}^T) = \det(\mathbf{A})^2$, so $|\det(\mathbf{A})| = \sqrt{\det(\mathbf{C})}$. Substituting everything back:
\begin{equation}
P(\mathbf{x}) = \frac{1}{\sqrt{(2\pi)^n \det(\mathbf{C})}} \exp\left( -\frac{1}{2} (\mathbf{x} - \boldsymbol{\mu})^T \mathbf{C}^{-1} (\mathbf{x} - \boldsymbol{\mu}) \right)
\end{equation}

When you move to the continuous limit ($n \to \infty$):
\begin{enumerate}
    \item The vector $\mathbf{x}$ becomes the path $X(t)$.
    \item The matrix multiplication $\mathbf{C}^{-1} \mathbf{x}$ becomes an integral $\int C^{-1}(t, t') X(t') dt'$.
    \item The quadratic form $(\mathbf{x}-\boldsymbol{\mu})^T \mathbf{C}^{-1} (\mathbf{x}-\boldsymbol{\mu})$ becomes the double integral:
\end{enumerate}
\begin{equation}
\int_0^T dt \int_0^T dt' (X(t) - \mu(t)) C^{-1}(t, t') (X(t') - \mu(t'))
\end{equation}
Using this we arrive at the final expression for the probability density functional of a Gaussian process:
\begin{equation}P[X] = \frac{1}{\sqrt{\det(2\pi C)}} \ \exp\left( -\frac{1}{2} \int_0^T dt \int_0^T dt' (X(t) - \mu(t)) C^{-1}(t, t') (X(t') - \mu(t')) \right).\end{equation}

\subsection{Onsager-Machlup Functional}

The \textbf{Onsager-Machlup functional} is essentially the ``Lagrangian'' for a stochastic process. Just as classical mechanics uses a Lagrangian to find the path of least action, stochastic mechanics uses this functional to find the \textbf{most probable path} a system will take.

Consider a particle governed by a Langevin equation with a deterministic force $f(x)$ (such as a gradient $-\nabla V$) and Gaussian white noise $\eta(t)$:
\begin{equation}
\frac{dx}{dt} = f(x) + \sqrt{2D} \eta(t)
\end{equation}
Here, $D$ is the diffusion coefficient. We already know from our previous derivation that the probability of a noise realization $\eta(t)$ is:
\begin{equation}
P[\eta] \propto \exp\left( -\frac{1}{4D} \int_0^T \eta(t)^2 \, dt \right)
\end{equation}
This equation is nothing but the continuous version of the joint distribution of the noise at different time points, analogous to $P(\mathbf{z})$ for the independent Gaussian variables.

To find the probability of a specific path $x(t)$, we substitute the expression for $\eta(t)$ from the Langevin equation:
\begin{equation}
\eta(t) = \frac{1}{\sqrt{2D}} \left( \dot{x} - f(x) \right)
\end{equation}
When we plug this into the noise probability functional, we get the core of the Onsager-Machlup action:
\begin{equation}
P[x] \propto \mathcal{J} \exp\left( -\frac{1}{4D} \int_0^T (\dot{x} - f(x))^2 \, dt \right)
\end{equation}

Because we changed variables from a function $\eta(t)$ to a function $x(t)$, we must include a \textbf{functional Jacobian}. In stochastic calculus, this term depends on the discretization (Itô vs. Stratonovich). For the standard Itô interpretation, the result is:
\begin{equation}
L(x, \dot{x}) = \frac{1}{4D} (\dot{x} - f(x))^2 + \frac{1}{2} \nabla f(x)
\end{equation}
The final probability functional for the path is:
\begin{equation}
P[x] \propto \exp\left( -\int_0^T \left[ \frac{(\dot{x} - f(x))^2}{4D} + \frac{1}{2} \nabla f(x) \right] dt \right)
\end{equation}

\begin{itemize}
    \item \textbf{The Most Probable Path:} If you want to know the most likely way a protein folds or a chemical reaction occurs, you minimize the integral of this functional (the ``Action'').
    \item \textbf{Transition Paths:} It allows us to calculate the probability of rare events, like a particle jumping over a high energy barrier.
    \item \textbf{Connection to Physics:} The term $(\dot{x} - f(x))^2$ is physically the ``extra'' work or dissipation required to deviate from the natural force $f(x)$.
\end{itemize}

\subsection{Path integral for Gaussian processes}
The path integral formulation for this Gaussian process can be derived by considering the probability of a particular path $X(t)$ given the noise $\eta(t)$ and integrating over all possible noise realizations. This leads to the expression for the probability density functional in terms of the action functional $S[X]$:
\begin{equation}P[X] \propto \exp\left( -S[X] \right)\end{equation}
where the action functional $S[X]$ is given by:
\begin{equation}S[X] = \frac{1}{2} \int_0^T dt \int_0^T dt' (X(t) - \mu(t)) C^{-1}(t, t') (X(t') - \mu(t'))\end{equation}
This action functional captures the dynamics of the Gaussian process and serves as the basis for computing various quantities of interest, such as transition probabilities, correlation functions, and response functions, by integrating over the space of all possible paths $X(t)$ weighted by $e^{-S[X]}$.

The probability of a stochastic trajectory $x(t)$ in the interval $[0, T]$ is determined by the \textbf{Action Functional} $S[x]$. The relationship between the Action and the \textbf{Onsager-Machlup Lagrangian} $L(x, \dot{x})$ is given by the integral:
\begin{equation}
S[x] = \int_0^T L(x(t), \dot{x}(t)) \, dt
\end{equation}
For a system governed by $\dot{x} = f(x) + \sqrt{2D}\eta(t)$, the explicit form is:
\begin{equation}
S[x] = \int_0^T \left[ \frac{(\dot{x} - f(x))^2}{4D} + \frac{1}{2} \nabla f(x) \right] dt
\end{equation}
The path $x^*(t)$ that minimizes this Action:
\begin{equation}
\frac{\delta S}{\delta x} = 0.
\end{equation}
defines the \textbf{Most Probable Path} of the system. This path represents the trajectory that requires the ``least effort'' from the thermal noise to deviate the particle from its deterministic force field. The minimization of the action results in Euler-Lagrange equations that govern the dynamics of the most probable path and is given by:
\begin{equation}\frac{d}{dt} \left( \frac{\partial L}{\partial \dot{x}} \right) - \frac{\partial L}{\partial x} = 0.\end{equation}
This framework allows us to analyze the behavior of stochastic systems, compute transition probabilities, and understand the dynamics of rare events by studying the most probable paths that the system can take under the influence of noise.

\subsubsection*{Example: Harmonic Potential}

To derive the probability distribution for a particle in a harmonic potential $V(X) = \frac{1}{2}kX^2$ from scratch, we must solve the path integral for the \textbf{Ornstein-Uhlenbeck process}. 

Because the Action is quadratic in $X$ and $\dot{X}$, the path integral is dominated by the \textbf{classical path} (the most probable path) that minimizes the action.
For a force $f(X) = -kX$, the Onsager-Machlup Lagrangian is:
\begin{equation}
L(X, \dot{X}) = \frac{(\dot{X} + kX)^2}{4D}
\end{equation}
(Note: We ignore the constant Jacobian term $\frac{k}{2}$ for now as it affects only the normalization $\mathcal{N}$). The Action $S[X]$ is:
\begin{equation}
S[X] = \int_0^T \frac{(\dot{X} + kX)^2}{4D} dt = \frac{1}{4D} \int_0^T (\dot{X}^2 + 2kX\dot{X} + k^2 X^2) dt
\end{equation}

The term $2kX\dot{X}$ can be written as a total derivative:
\begin{equation}
2kX\frac{dX}{dt} = k \frac{d}{dt}(X^2)
\end{equation}
Integrating this over $[0, T]$ gives:
\begin{equation}
\int_0^T k \frac{d}{dt}(X^2) dt = k(X_T^2 - X_0^2)
\end{equation}
The Action becomes:
\begin{equation}
S[X] = \frac{k(X_T^2 - X_0^2)}{4D} + \frac{1}{4D} \int_0^T (\dot{X}^2 + k^2 X^2) dt
\end{equation}

The path that minimizes the integral $\int (\dot{X}^2 + k^2 X^2) dt$ must satisfy:
\begin{equation}
\frac{d}{dt}\left(\frac{\partial L}{\partial \dot{X}}\right) - \frac{\partial L}{\partial X} = 0 \implies \ddot{X} - k^2 X = 0
\end{equation}
The general solution is:
\begin{equation}
X^*(t) = A e^{kt} + B e^{-kt}
\end{equation}
Using boundary conditions $X(0) = X_0$ and $X(T) = X_T$:
\begin{itemize}
    \item $A + B = X_0$
    \item $A e^{kT} + B e^{-kT} = X_T$
\end{itemize}
Solving for $A$ and $B$:
\begin{equation}
A = \frac{X_T - X_0 e^{-kT}}{2 \sinh(kT)}, \quad B = \frac{X_0 e^{kT} - X_T}{2 \sinh(kT)}
\end{equation}

Substituting $X^*(t)$ back into the Action $S[X^*]$ (after some heavy algebraic simplification of the hyperbolic functions):
\begin{equation}
S[X^*] = \frac{k}{2D} \left[ \frac{(X_T - X_0 e^{-kT})^2}{1 - e^{-2kT}} \right]
\end{equation}
The transition probability is $P(X_T, T | X_0, 0) \propto \exp(-S[X^*])$. Including the normalization factor $\mathcal{N}$:
\begin{equation}
P(X_T, T | X_0, 0) = \sqrt{\frac{k}{2\pi D(1 - e^{-2kT})}} \exp\left( -\frac{k(X_T - X_0 e^{-kT})^2}{2D(1 - e^{-2kT})} \right).
\end{equation}

A simple example when the potential is zero ($k=0$) gives us the free particle case:
\begin{equation}P(X_T, T | X_0, 0) = \frac{1}{\sqrt{4\pi D T}} \exp\left( -\frac{(X_T - X_0)^2}{4D T} \right)\end{equation}
This is the well-known result for a free particle undergoing Brownian motion, where the transition probability is given by a Gaussian distribution. This illustrates how the path integral formulation can be used to derive transition probabilities for stochastic processes, and how it can be applied to more complex systems by incorporating potential functions and interactions.

\subsection{Fokker-Planck equation from Path Integral}

To arrive at the \textbf{Fokker-Planck equation} from the path integral $P[X] \propto \exp(-S[X])$, we use the \textbf{Schr\"{o}dinger-like evolution} of the probability density. The transition probability $P(x, t+\Delta t | x', t)$ acts as a propagator that moves the distribution forward by an infinitesimal time step $\Delta t$.

The probability density $P(x, t+\Delta t)$ is related to the density at the previous time step $P(x', t)$ by integrating over all possible starting positions $x'$ (also known as the Chapman-Kolmogorov equation):
\begin{equation}
P(x, t+\Delta t) = \int_{-\infty}^{\infty} P(x, t+\Delta t | x', t) P(x', t) dx'
\end{equation}
From the path integral definition, the transition probability for a small time $\Delta t$ is determined by the Action $S[X] = \int L dt$. Using the Onsager-Machlup Lagrangian $L = \frac{(\dot{x}-f(x))^2}{4D}$, and approximating $\dot{x} \approx \frac{x-x'}{\Delta t}$:
\begin{equation}
P(x, t+\Delta t | x', t) = \frac{1}{\sqrt{4\pi D \Delta t}} \exp\left( -\frac{(x - x' - f(x')\Delta t)^2}{4D\Delta t} \right)
\end{equation}
Let $\xi = x - x'$ be the displacement. Then $x' = x - \xi$. We rewrite the integral in terms of $\xi$:
\begin{equation}
P(x, t+\Delta t) = \int_{-\infty}^{\infty} \frac{1}{\sqrt{4\pi D \Delta t}} \exp\left( -\frac{(\xi - f(x-\xi)\Delta t)^2}{4D\Delta t} \right) P(x-\xi, t) d\xi
\end{equation}
We expand $P(x-\xi, t)$ and $f(x-\xi)$ around $x$. Using the Gaussian moments for $\langle \xi \rangle \approx f(x)\Delta t$ and $\langle \xi^2 \rangle \approx 2D\Delta t$:
\begin{equation}
P(x-\xi, t) \approx P(x, t) - \xi \frac{\partial P}{\partial x} + \frac{\xi^2}{2} \frac{\partial^2 P}{\partial x^2}
\end{equation}
\begin{equation}
f(x-\xi) \approx f(x) - \xi \frac{\partial f}{\partial x}
\end{equation}
% Substitute the expansions back into the integral and collect terms of order $\Delta t$:
% \begin{equation}
% P(x, t+\Delta t) \approx P(x, t) - \Delta t \frac{\partial}{\partial x} [f(x)P(x, t)] + D\Delta t \frac{\partial^2 P(x, t)}{\partial x^2}
% \end{equation}
Substituting the Taylor expansions into the Chapman-Kolmogorov integral:
\begin{equation}
P(x, t+\Delta t) = \int G(\xi) \left[ P(x,t) - \xi \frac{\partial P}{\partial x} + \frac{\xi^2}{2} \frac{\partial^2 P}{\partial x^2} \right] d\xi
\end{equation}

We evaluate the integral term-by-term using Gaussian moments $\langle 1 \rangle = 1$, $\langle \xi \rangle = f(x)\Delta t$, and $\langle \xi^2 \rangle = 2D\Delta t$:
\begin{itemize}
    \item \textbf{Term A:} $\int G(\xi) P(x,t) d\xi = P(x,t)$
    \item \textbf{Term B:} $-\int \xi G(\xi) \frac{\partial P}{\partial x} d\xi = -f(x)\Delta t \frac{\partial P}{\partial x}$
    \item \textbf{Term C:} $\int \frac{\xi^2}{2} G(\xi) \frac{\partial^2 P}{\partial x^2} d\xi = D\Delta t \frac{\partial^2 P}{\partial x^2}$
\end{itemize}
Summing these results gives the evolution to order $\Delta t$:
\begin{equation}
P(x, t+\Delta t) \approx P(x, t) - \Delta t \left( f(x) \frac{\partial P}{\partial x} + P \frac{\partial f}{\partial x} \right) + D\Delta t \frac{\partial^2 P}{\partial x^2}
\end{equation}
Recognizing the product rule $\frac{\partial}{\partial x}(fP) = f \partial_x P + P \partial_x f$, we obtain:
\begin{equation}
P(x, t+\Delta t) - P(x, t) = \Delta t \left[ -\frac{\partial}{\partial x} (f P) + D \frac{\partial^2 P}{\partial x^2} \right]
\end{equation}
Dividing by $\Delta t$ and taking the limit $\Delta t \to 0$ yields the Fokker-Planck equation:
\begin{equation}
\frac{\partial P}{\partial t} = -\frac{\partial}{\partial x} [f(x)P] + D \frac{\partial^2 P}{\partial x^2}
\end{equation}

\end{singlespace}
\end{document}